\begin{theorem}[Functional equation for Hecke L-functions] \label{theorem:functional_equation_for_a_hecke_L_function_of_a_hecke_character_of_a_number_field}
\TODO{TODO: define the conductor of a hecke character}
Let $K$ be a number field with absolute discriminant $|\Delta_{K/\mathbb{Q}}|$, and let $\chi$ be a \CrefAndHyperrefIfExist{definition:hecke_character_of_a_number_field}{Hecke character of $K$} with conductor $\mathfrak{f}(\chi)$. The \CrefAndHyperrefIfExist{definition:completed_hecke_L_function_of_a_hecke_character_of_a_number_field}{completed Hecke $L$-function} $\Lambda(s,\chi)$ satisfies the function equation
\[ \boxed{ \Lambda(s, \chi) = \varepsilon(\chi) \, \Lambda(1 - s, \overline{\chi}), } \]
\TODO{TODO: define root number}
\TODO{Argue that the Hecke $L$-function is meromorphic; this function equation should say that the $L$-function is meromorphic for the real part of $s$ greater than $1$ or less than $0$, but an argument for what happens in the critical strip is necessary.}
where $\varepsilon(\chi) \in \mathbb{C}$ is the root number with $|\varepsilon(\chi)| = 1$, and $\overline{\chi}$ is the complex conjugate Hecke character.
\end{theorem}